% Lean: EconLean/Foundations/FixedPoint.lean (vendored proofs in Foundations/FixedPoint/),
%       EconLean/GE/ExistenceGeneral.lean
\chapter{Fixed-Point Theorems and General Existence}
\label{chap:fixedpoint}

The proofs of Brouwer's and Kakutani's theorems are vendored from \texttt{vltanh/lean4-fixed-point} (The-Anh Vu-Le,
Apache 2.0), which proves them through Sperner's lemma and the Knaster--Kuratowski--Mazurkiewicz (KKM) lemma. This
chapter records the economics-facing statements and the step that makes equilibrium existence unconditional.

\begin{theorem}[Brouwer's fixed-point theorem]
\label{thm:brouwer}
\lean{Econ.brouwer_fixedPoint, Econ.brouwer_fixedPoint_continuousMap}
\leanok
Let $s$ be a nonempty compact convex subset of a finite-dimensional real normed space and $f$ continuous on $s$ with
$f(s) \subseteq s$. Then $f(x) = x$ for some $x \in s$.
\end{theorem}
\begin{proof}
\leanok
This is \texttt{Brouwer.KKM.exists\_fixedPoint} (vendored). Its proof: (1) Sperner's lemma gives, in each barycentric
subdivision of the standard simplex $\Delta$, a fully labelled cell; (2) the KKM lemma follows by compactness, as the
cells shrink; (3) applying KKM to $F_i = \{x \in \Delta : f(x)_i \le x_i\}$ gives a point with $f(x)_i \le x_i$ for all
$i$, and both sum to $1$, so $f(x) = x$; (4) a compact convex set $s$ lies in a simplex, and composing $f$ with the
nearest-point retraction onto $s$ and barycentric coordinates transfers the fixed point to $s$; (5) a
finite-dimensional normed space is linearly homeomorphic to a Euclidean space.
\end{proof}

\begin{theorem}[Kakutani's fixed-point theorem]
\label{thm:kakutani}
\lean{Econ.kakutani_fixedPoint}
\leanok
\uses{thm:brouwer}
Let $s$ be a nonempty compact convex subset of a finite-dimensional real normed space and $\varphi$ assign to each
$x \in s$ a nonempty convex $\varphi(x) \subseteq s$, with closed graph $\{(x, y) : x \in s,\ y \in \varphi(x)\}$. Then
$x \in \varphi(x)$ for some $x \in s$.
\end{theorem}
\begin{proof}
\leanok
This is \texttt{Kakutani.exists\_fixedPoint} (vendored), Kakutani's 1941 proof. On the standard simplex: pick
$y_v \in \varphi(v)$ for each $v$; the map that is linear on each cell of the $n$-th barycentric subdivision with value
$y_v$ at each vertex $v$ is continuous, so Brouwer gives a fixed point $x_n = \sum_i \lambda_{ni}\, y(x_{ni})$, where
the $x_{ni}$ are the vertices of its cell, within $2(r/(r+1))^n$ of $x_n$. Along a subsequence $x_n \to x_0$,
$\lambda_n \to \lambda_0$ and $y(x_{ni}) \to y_{0i}$; the closed graph gives $y_{0i} \in \varphi(x_0)$, and convexity
gives $x_0 = \sum_i \lambda_{0i} y_{0i} \in \varphi(x_0)$. For general $s$: $s$ lies in a simplex $K$; apply the simplex
case to $x \mapsto \varphi(\pi(x))$ on $K$, with $\pi$ the nearest-point retraction onto $s$; the fixed point lies in
$s$, where $\pi$ is the identity. A linear isomorphism with Euclidean space handles normed spaces.
\end{proof}

\begin{lemma}[The price simplex is a Brouwer domain]
\label{lem:priceSimplex_brouwer}
\lean{Econ.convex_priceSimplex, Econ.isCompact_priceSimplex, Econ.priceSimplex_nonempty, Econ.brouwerSimplex}
\leanok
\uses{def:priceSimplex, thm:brouwer}
$\Delta$ is convex, compact, and nonempty when there is at least one good; hence Brouwer's theorem holds on $\Delta$.
\end{lemma}
\begin{proof}
\leanok
Convexity: nonnegativity and $\sum_i p_i = 1$ are preserved by convex combinations. Compactness: $\Delta$ is the
intersection of the closed sets $\{p \ge 0\}$ and $\{\sum_i p_i = 1\}$ and lies in the compact cube $[0,1]^\iota$
(Lemma: coordinates are at most $1$). Nonemptiness: uniform prices $p_i = 1/|\iota|$. Then apply
Theorem~\ref{thm:brouwer} in $\R^\iota$.
\end{proof}

\begin{theorem}[Unconditional existence]
\label{thm:existence_unconditional}
\lean{Econ.gale_nikaido_debreu', Econ.ExchangeEconomy.exists_freeDisposalEquilibrium'}
\leanok
\uses{lem:priceSimplex_brouwer, thm:gnd, thm:ge_existence}
For any nonempty finite set of goods, the Gale--Nikaido--Debreu lemma and the existence of a free-disposal equilibrium
(Theorems~\ref{thm:gnd}, \ref{thm:ge_existence}) hold without the Brouwer hypothesis. Demand must still be defined and
continuous on all of $\Delta$, so strongly monotone economies (MWG Prop.~17.C.1) are not yet covered.
\end{theorem}
\begin{proof}
\leanok
Supply Lemma~\ref{lem:priceSimplex_brouwer} as the Brouwer hypothesis of the conditional theorems.
\end{proof}

\section{Existence with a boundary condition}
\label{sec:existence_boundary}

\begin{definition}[Bounded below; boundary condition; MWG Prop.~17.B.2(iv),(v)]
\label{def:boundary_condition}
\lean{Econ.ExcessBoundedBelow, Econ.BoundaryCondition, Econ.PositivePrices}
\leanok
$z$ is \emph{bounded below} if there is $s$ with $z_l(p) > -s$ for all $p \gg 0$ and all $l$. $z$ satisfies the
\emph{boundary condition} if whenever $p^n \gg 0$, $p^n \to \bar p \ne 0$ and $\bar p_l = 0$ for some $l$, then for every
$M$, eventually $z_l(p^n) > M$ for some $l$ (that is, $\max_l z_l(p^n) \to \infty$).
\end{definition}

\begin{lemma}[Shrinking the simplex]
\label{lem:shrink_price}
\lean{Econ.shrinkPrice, Econ.isOpen_positivePrices, Econ.continuous_shrinkPrice, Econ.shrinkPrice_mem,
  Econ.cost_shrinkPrice}
\leanok
\uses{def:priceSimplex}
Let $N$ be the number of goods and $r_\varepsilon(p) = (1 - N\varepsilon)p + \varepsilon \mathbf 1$. For $0 < \varepsilon$,
$N\varepsilon \le 1$, $r_\varepsilon$ is continuous and maps $\Delta$ into strictly positive prices in $\Delta$, and
$r_\varepsilon(p) \cdot x = (1 - N\varepsilon)\, p \cdot x + \varepsilon \sum_k x_k$. The set of strictly positive prices
is open.
\end{lemma}
\begin{proof}
\leanok
Coordinates are $(1 - N\varepsilon)p_i + \varepsilon \ge \varepsilon > 0$, and they sum to $(1 - N\varepsilon) + N\varepsilon = 1$.
The cost identity is linearity. Positive prices are a finite intersection of open half-spaces.
\end{proof}

\begin{lemma}[Approximate equilibria]
\label{lem:approx_equilibrium}
\lean{Econ.exists_approx_equilibrium}
\leanok
\uses{lem:shrink_price, thm:existence_unconditional}
If $z$ is continuous on $p \gg 0$ with Walras' law there, then for $0 < \varepsilon$, $N\varepsilon \le 1$ there is
$x \gg 0$ in $\Delta$ with $(1 - N\varepsilon) z_l(x) + \varepsilon \sum_k z_k(x) \le 0$ for every $l$.
\end{lemma}
\begin{proof}
\leanok
Let $w(p)_l = (1 - N\varepsilon) z_l(r_\varepsilon p) + \varepsilon \sum_k z_k(r_\varepsilon p)$, continuous on $\Delta$. Since
$\sum_i p_i = 1$, $p \cdot w(p) = (1 - N\varepsilon)\, p \cdot z(r_\varepsilon p) + \varepsilon \sum_k z_k(r_\varepsilon p)
= r_\varepsilon(p) \cdot z(r_\varepsilon p) = 0$. The unconditional Gale--Nikaido--Debreu lemma gives $p \in \Delta$ with
$w(p) \le 0$; take $x = r_\varepsilon(p)$.
\end{proof}

\begin{theorem}[Existence of equilibrium; MWG Prop.~17.C.1]
\label{thm:existence_boundary}
\lean{Econ.exists_equilibrium_of_boundaryCondition}
\leanok
\uses{def:boundary_condition, lem:approx_equilibrium, lem:priceSimplex_brouwer}
If $z$ is continuous on $p \gg 0$, satisfies Walras' law there, is bounded below and satisfies the boundary condition,
then $z(p) = 0$ for some $p \gg 0$. (Homogeneity is not needed.)
\end{theorem}
\begin{proof}
\leanok
Replace $s$ by $S = \max(s, 0)$. Let $\varepsilon_n = \frac{1}{2N(n+1)}$, so $N\varepsilon_n \le \frac12$ and
$\varepsilon_n \to 0$, and let $x_n$ be the approximate equilibria of Lemma~\ref{lem:approx_equilibrium}.

\emph{Upper bound.} If $z_l(x_n) > S$, then since $\sum_k z_k(x_n) > -NS$,
$0 \ge (1 - N\varepsilon_n) z_l(x_n) + \varepsilon_n \sum_k z_k(x_n) > \tfrac12 S - N\varepsilon_n S \ge 0$, a contradiction. So
$z_l(x_n) \le S$ for all $n, l$.

\emph{Limit.} $\Delta$ is compact, so a subsequence $x_{\varphi(n)} \to \bar p \in \Delta$, and $\bar p \ne 0$. If
$\bar p_l = 0$ for some $l$, the boundary condition with $M = S$ gives $z_k(x_{\varphi(n)}) > S$ for some $n, k$,
contradicting the upper bound. So $\bar p \gg 0$, $z$ is continuous at $\bar p$, and letting $n \to \infty$ in
$(1 - N\varepsilon) z_l(x) + \varepsilon \sum_k z_k(x) \le 0$ gives $z_l(\bar p) \le 0$ for every $l$.

\emph{Equilibrium.} Walras' law $\sum_i \bar p_i z_i(\bar p) = 0$ is a sum of nonpositive terms, so each is $0$, and
$\bar p_i > 0$ gives $z_i(\bar p) = 0$.
\end{proof}

\begin{theorem}[Boundary-condition example]
\label{thm:zcd_boundary}
\lean{Econ.AggregationExamples.zCD_continuousOn, Econ.AggregationExamples.zCD_boundedBelow,
  Econ.AggregationExamples.eventually_lt_div_sub_one, Econ.AggregationExamples.zCD_boundary}
\leanok
\uses{thm:existence_boundary, thm:zcd_example}
$z(p) = (p_1/p_0 - 1, p_0/p_1 - 1)$ is continuous on $p \gg 0$, bounded below by $-1$, and satisfies the boundary
condition, so Theorem~\ref{thm:existence_boundary} applies. It is undefined (blows up) at the corners of $\Delta$, so the
closed-simplex theorem does not.
\end{theorem}
\begin{proof}
\leanok
Continuity: quotients of coordinates with nonzero denominators. Lower bound: $p_1/p_0 > 0$. Boundary: if
$p^n \to \bar p \ne 0$ with $\bar p_l = 0$, then $\bar p \ge 0$ forces the other coordinate $\bar p_k > 0$; as
$p^n_l \to 0^+$, $1/p^n_l \to \infty$, and $p^n_k \to \bar p_k > 0$, so $z_l(p^n) = p^n_k/p^n_l - 1 \to \infty$.
\end{proof}

\begin{theorem}[Existence in an exchange economy, excess-demand form]
\label{thm:existence_boundary_economy}
\lean{Econ.ExchangeEconomy.exists_walrasianEquilibrium_of_boundaryCondition}
\leanok
\uses{thm:existence_boundary, lem:agg_walras}
With locally nonsatiated consumers and a demand selection on $p \gg 0$ whose aggregate excess demand is continuous there,
bounded below and satisfies the boundary condition, a Walrasian equilibrium with $p \gg 0$ exists.
\end{theorem}
\begin{proof}
\leanok
Walras' law for aggregate excess demand follows from local nonsatiation (Lemma~\ref{lem:agg_walras}). Theorem~\ref{thm:existence_boundary}
gives $p \gg 0$ with zero excess demand, which is a Walrasian equilibrium.
\end{proof}

\section{Existence from preferences}
\label{sec:existence_preferences}

\begin{lemma}[Boundary condition from strong monotonicity; MWG Prop.~17.B.2(v)]
\label{lem:boundary_from_prefs}
\lean{Econ.ExchangeEconomy.boundaryCondition_excessDemand}
\leanok
\uses{def:boundary_condition}
If utilities are continuous, preferences strongly monotone and $\bar\omega \gg 0$, aggregate excess demand of any demand
selection on $p \gg 0$ satisfies the boundary condition.
\end{lemma}
\begin{proof}
\leanok
Let $p^n \gg 0$, $p^n \to \bar p \ne 0$, $\bar p_l = 0$. Then $\bar p \ge 0$, and $\bar p \cdot \bar\omega > 0$, so some consumer
$a$ has $\bar p \cdot \omega_a > 0$.

\emph{Claim: for every $M'$, eventually $d_{a,k}(p^n) > M'$ for some $k$.} Otherwise, frequently $d_a(p^n) \le M'$
coordinatewise; extract a subsequence along which $d_a(p^n)$ lies in the box $[0, M']^\iota$, and a further one along which
$x_n = d_a(p^n) \to \bar x \ge 0$. Passing to the limit in $p^n \cdot x_n \le p^n \cdot \omega_a$ gives
$\bar p \cdot \bar x \le \bar p \cdot \omega_a$. Let $e = e_l$; since $\bar p_l = 0$, $\bar p \cdot e = 0$, and by strong monotonicity
$u_a(\bar x + e) > u_a(\bar x)$. By continuity in $t$ there is $t \in [0, 1)$ with $u_a(t\bar x + e) > u_a(\bar x)$. Then
$y = t\bar x + e \ge 0$ and $\bar p \cdot y = t\, \bar p \cdot \bar x \le t\, \bar p \cdot \omega_a < \bar p \cdot \omega_a$. By continuity
of costs, $p^n \cdot y < p^n \cdot \omega_a$ for large $n$, so $y$ is affordable and $u_a(x_n) \ge u_a(y)$. In the limit
$u_a(\bar x) \ge u_a(y) > u_a(\bar x)$, a contradiction.

\emph{Aggregate.} Demands are nonnegative, so $z_k(p^n) \ge d_{a,k}(p^n) - \bar\omega_k \ge d_{a,k}(p^n) - \sum_j \bar\omega_j$.
Apply the claim with $M' = M + \sum_j \bar\omega_j$.
\end{proof}

\begin{theorem}[Existence from preferences; MWG Props.~17.B.2 and 17.C.1]
\label{thm:existence_preferences}
\lean{Econ.ExchangeEconomy.wealth_pos, Econ.ExchangeEconomy.exists_walrasianEquilibrium_of_preferences}
\leanok
\uses{lem:boundary_from_prefs, thm:existence_boundary_economy}
In an exchange economy where each consumer has a continuous, strictly quasiconcave utility representing strongly
monotone preferences and a nonzero endowment, and $\bar\omega \gg 0$, a Walrasian equilibrium with $p \gg 0$ exists.
\end{theorem}
\begin{proof}
\leanok
At $p \gg 0$ each consumer's wealth $p \cdot \omega_a$ is positive (nonzero, nonnegative endowment). Strict quasiconcavity
gives a unique demand, which is continuous in $(p, w)$ by Berge's theorem; composing with $p \mapsto (p, p \cdot \omega_a)$
gives a demand selection continuous on $p \gg 0$, so excess demand is continuous there. It is bounded below by
$-\sum_j \bar\omega_j - 1$ since demands are nonnegative. Strong monotonicity gives local nonsatiation and, by
Lemma~\ref{lem:boundary_from_prefs}, the boundary condition. Apply Theorem~\ref{thm:existence_boundary_economy}.
\end{proof}

\begin{theorem}[Square-root economy]
\label{thm:sqrt_economy}
\lean{Econ.GEExamples.uSqrt, Econ.GEExamples.continuous_uSqrt, Econ.GEExamples.strictlyQuasiconcave_uSqrt,
  Econ.GEExamples.stronglyMonotone_uSqrt, Econ.GEExamples.sqrtEconomy,
  Econ.GEExamples.sqrtEconomy_exists_equilibrium}
\leanok
\uses{thm:existence_preferences}
With $u(x) = \sqrt{x_0} + \sqrt{x_1}$ for both consumers and endowments $(1,2), (2,1)$, a Walrasian equilibrium with
strictly positive prices exists. Demand does not exist at prices with a zero coordinate, so the closed-simplex theorem
does not apply.
\end{theorem}
\begin{proof}
\leanok
$\sqrt{\cdot}$ is continuous and strictly concave on $[0, \infty)$ and strictly increasing there. Strict quasiconcavity:
$t u(x) + (1 - t) u(y) \le u(tx + (1-t)y)$ coordinatewise by concavity, strictly in a coordinate where $x, y$ differ, and
$\min(u(x), u(y)) \le t u(x) + (1-t) u(y)$. Strong monotonicity: each square root weakly increases and one strictly. Apply
Theorem~\ref{thm:existence_preferences}.
\end{proof}

\section{Existence with zero endowments and Cobb--Douglas-type preferences}
\label{sec:existence_regular}

\begin{definition}[Regular and interior-regular utilities]
\label{def:regular_utility}
\lean{Econ.RegularUtility, Econ.InteriorRegular, Econ.RegularUtility.demand_eq_empty_of_boundary}
\leanok
A utility $u$ is \emph{regular} if it is continuous, demand is at most a singleton at $p \gg 0$ and $w > 0$, it is locally
nonsatiated, and it has the \emph{improvement property}: whenever $\bar x \ge 0$, $\bar p \ge 0$ has a zero coordinate, and
$\bar p \cdot \bar x \le W$ with $W > 0$, some $y \ge 0$ has $u(y) > u(\bar x)$ and $\bar p \cdot y < W$. It is
\emph{interior-regular} (Cobb--Douglas type) if it is continuous, strictly increasing from interior bundles, strictly
quasiconcave on the interior, and every interior bundle is strictly better than every boundary bundle.
\end{definition}

\begin{lemma}[Both classical hypothesis sets give regularity]
\label{lem:regular_from}
\lean{Econ.exists_lt_one_better, Econ.cost_smul_add_single_lt, Econ.single_nonneg',
  Econ.regularUtility_of_strictlyQuasiconcave, Econ.positive_add_const, Econ.InteriorRegular.locallyNonsatiated,
  Econ.regularUtility_of_interiorRegular, Econ.RegularUtility.exists_continuous_demand_function}
\leanok
\uses{def:regular_utility}
Continuous, strictly quasiconcave, strongly monotone utilities are regular. So are interior-regular utilities. A
regular utility has a continuous demand function on $p \gg 0$, $w > 0$.
\end{lemma}
\begin{proof}
\leanok
Improvement: if $u(\bar x + e_l) > u(\bar x)$ with $\bar p_l = 0$, continuity gives $t \in [0,1)$ with
$u(t\bar x + e_l) > u(\bar x)$. This $y$ costs $t\,\bar p\cdot\bar x < W$. Strong monotonicity gives the first inequality. In the
interior-regular case it holds when $\bar x \gg 0$; if $\bar x$ is on the boundary, use the cheap interior bundle
$y = \frac{W}{2\sum_i \bar p_i + 1}\mathbf 1$ instead. For uniqueness in the interior-regular case, every demanded bundle is
interior: a boundary bundle is beaten by the affordable interior bundle above. So two different demanded bundles would
have a strictly better affordable midpoint. Local nonsatiation: $x + \frac{\varepsilon}{2}\mathbf 1$ is interior and
strictly better. Continuity of demand is Berge's theorem, as in Theorem~\ref{thm:existence_preferences}.
\end{proof}

\begin{theorem}[Existence for regular consumers]
\label{thm:existence_regular}
\lean{Econ.ExchangeEconomy.demand_zero_wealth, Econ.ExchangeEconomy.boundaryCondition_excessDemand_of_regular,
  Econ.ExchangeEconomy.exists_walrasianEquilibrium_of_regular,
  Econ.ExchangeEconomy.exists_walrasianEquilibrium_of_preferences_of_totalEndowment,
  Econ.ExchangeEconomy.exists_walrasianEquilibrium_of_interiorRegular}
\leanok
\uses{lem:regular_from, lem:boundary_from_prefs, thm:existence_boundary_economy}
If every consumer's utility is regular and $\bar\omega \gg 0$, a Walrasian equilibrium with $p \gg 0$ exists, with no
assumption on individual endowments. In particular this holds for the hypotheses of Theorem~\ref{thm:existence_preferences}
without $\omega_a \ne 0$, and for interior-regular (Cobb--Douglas-type) consumers.
\end{theorem}
\begin{proof}
\leanok
A consumer with $\omega_a = 0$ has budget set $\{0\}$ at $p \gg 0$ and demands $0$, a constant, hence continuous, selection.
The others have positive wealth and a continuous demand function. The boundary condition follows as in
Lemma~\ref{lem:boundary_from_prefs}, using the improvement property in place of strong monotonicity. Apply
Theorem~\ref{thm:existence_boundary_economy}.
\end{proof}

\begin{theorem}[A Cobb--Douglas consumer and a consumer who owns nothing]
\label{thm:cd_economy}
\lean{Econ.GEExamples.uCD, Econ.GEExamples.interiorRegular_uCD, Econ.GEExamples.cdEconomy,
  Econ.GEExamples.cdEconomy_endow_one, Econ.GEExamples.cdEconomy_exists_equilibrium,
  Econ.GEExamples.cdEconomy_trade}
\leanok
\uses{thm:existence_regular, thm:sqrt_economy}
Consumer $0$ has $u(x) = x_0 x_1$ and owns $(1,0)$; consumers $1$ and $2$ have $u(x) = \sqrt{x_0} + \sqrt{x_1}$, and
own nothing and $(0,1)$ respectively. A Walrasian equilibrium with $p \gg 0$ exists, and in every such equilibrium
consumer $0$ trades: its endowment has utility $0$, while the affordable bundle $(1/2, p_0/(2p_1))$ has positive utility.
\end{theorem}
\begin{proof}
\leanok
$x_0x_1$ is interior-regular. For strict quasiconcavity on the interior, write $A = x_0x_1$, $B = y_0y_1$ and
$C = x_0y_1 + y_0x_1$. Then the mixture has utility $t^2A + s^2B + tsC$ with $s = 1 - t$, and
$C^2 - 4AB = (x_0y_1 - y_0x_1)^2 \ge 0$. If $A \le B$, this exceeds $A = (t+s)^2 A$ because $C \ge 2A$, with strict inequality
unless $A = B$ and $C = 2A$, which forces $x = y$. Apply Theorem~\ref{thm:existence_regular}.
\end{proof}

\begin{theorem}[Cobb--Douglas economies for every exponent]
\label{thm:cobb_douglas_existence}
\lean{Econ.cobbDouglas, Econ.logCobbDouglas, Econ.cobbDouglas_eq_exp, Econ.continuous_cobbDouglas,
  Econ.logCobbDouglas_strictMono, Econ.logCobbDouglas_strictConcave, Econ.positive_combo,
  Econ.interiorRegular_cobbDouglas, Econ.ExchangeEconomy.exists_walrasianEquilibrium_cobbDouglas,
  Econ.GEExamples.cd2Economy, Econ.GEExamples.cd2Economy_exists_equilibrium}
\leanok
\uses{def:regular_utility, lem:regular_from, thm:existence_regular}
For any exponents $\alpha_i > 0$, $u(x) = \prod_i x_i^{\alpha_i}$ is interior-regular. Hence every exchange economy of
Cobb--Douglas consumers with $\bar\omega \gg 0$ has a Walrasian equilibrium with $p \gg 0$, even if some consumers own
nothing. Witness: exponents $(1/3, 2/3)$ owning $(1,0)$ and $(3,1)$ owning $(0,1)$.
\end{theorem}
\begin{proof}
\leanok
On the interior $u = \exp(\sum_i \alpha_i \log x_i)$. The sum is strictly increasing, and strictly concave because $\log$
is and $x \ne y$ differ in some coordinate; $\exp$ preserves both. On the boundary some factor is $0^{\alpha_i} = 0$, while
interior utility is positive. Continuity: $x \mapsto x^{\alpha}$ is continuous on $\mathbb{R}$ for $\alpha > 0$.
\end{proof}

\begin{theorem}[Cobb--Douglas demand and an explicit equilibrium]
\label{thm:cobb_douglas_demand}
\lean{Econ.cdDemand, Econ.cost_cdDemand, Econ.demand_cobbDouglas, Econ.GEExamples.cd2Economy_equilibrium}
\leanok
\uses{thm:cobb_douglas_existence}
For $\alpha \gg 0$, $p \gg 0$ and $w > 0$, the demand of $u(x) = \prod_i x_i^{\alpha_i}$ is the single bundle
$x_i = \alpha_i w / ((\sum_j \alpha_j) p_i)$ (MWG Example 3.D.1). In the economy of Theorem~\ref{thm:cobb_douglas_existence},
prices $(1, 8/9)$ with allocations $(1/3, 3/4)$ and $(2/3, 1/4)$ form a Walrasian equilibrium.
\end{theorem}
\begin{proof}
\leanok
The bundle exhausts the budget. For any affordable $y$: if some $y_i = 0$ then $u(y) = 0 < u(x)$; otherwise
$\sum_i \alpha_i \log(y_i/x_i) \le \sum_i \alpha_i (y_i/x_i - 1) = \frac{\sum_j \alpha_j}{w} p \cdot y - \sum_j \alpha_j \le 0$.
Uniqueness follows from strict quasiconcavity on the interior. The equilibrium is a direct computation with the formula.
\end{proof}
